% =====================================================================
\section[Model 1 --- Estimation and Geometric Control]{Model 1 --- Nonlinear Model-Based Estimation and Geometric Control}
\label{sec:model1}
% =====================================================================

Model~1 is the loop that actually flies the drone. It starts from physics (block~1). It builds the
small-error linear model that an estimator needs (blocks 2--3) and estimates the state from the sensors (block~4).
Finally it computes the commands in two nested loops, position (blocks 5--6) and attitude (block~7), and turns them
into motor speeds (block~8).

% ---------------------------------------------------------------------
\subsection{Block 1 --- Physical system and nonlinear dynamics}
\label{sec:m1b1}
% ---------------------------------------------------------------------

The drone is modelled as a single rigid body with mass $m$ and inertia $J$, pushed by a thrust $T$ along its own
$z$-axis and twisted by a torque $\vtau$. Its state is position, velocity, attitude and angular velocity.

\begin{eqbox}{cPlant}{Equations M1.1 -- M1.4 \quad Rigid-body dynamics of the quadrotor}
\begin{align}
  \dot{\vp} &= \vv \tag{M1.1}\label{eq:m1.1}\\
  \dot{\vv} &= \vg + \frac{T}{m}\,R\,\ve_3 \tag{M1.2}\label{eq:m1.2}\\
  \dot{R}   &= R\,\skw{\vw} \tag{M1.3}\label{eq:m1.3}\\
  \dot{\vw} &= J^{-1}\big(\vtau - \vw\times J\vw\big) \tag{M1.4}\label{eq:m1.4}
\end{align}
Compactly $\dot{\vx}=\vf(\vx,\vu)$ with state $\vx=(\vp,\vv,R,\vw)$ and input $\vu=[T;\vtau]$.
\end{eqbox}

\begin{dimtable}
\vp,\ \dot{\vp} & position and its rate (world frame) & \R^3 & m, m/s\\
\vv,\ \dot{\vv} & velocity and acceleration (world frame) & \R^3 & m/s, m/s$^2$\\
R,\ \dot R & attitude and its rate & \R^{3\times3} & --, 1/s\\
\vw,\ \dot{\vw} & body angular velocity and acceleration & \R^3 & rad/s, rad/s$^2$\\
T & total thrust & \R & N\\
\vtau & body torque & \R^3 & N\,m\\
\ve_3 & unit vector $[0,0,1]^\top$ & \R^3 & --\\
\vx & state (stored as 18 numbers, 12 degrees of freedom) & \R^3{\times}\R^3{\times}\SO{\times}\R^3 & mixed\\
\vu & input & \R^4 & N, N\,m\\
\end{dimtable}

\subsubsection*{Derivation}

\textbf{(M1.1) Kinematics.} Velocity is by definition the time derivative of position: $\vv=\frac{d\vp}{dt}$.

\textbf{(M1.2) Newton's second law.} In the inertial (world) frame, $m\dot{\vv}=\sum\vF$. Two forces act:
\begin{itemize}
  \item gravity, $m\vg$;
  \item rotor thrust. In the body frame it is $T\ve_3$ (straight ``up'' out of the propeller discs). Rotating it into
        the world frame gives $R\,T\ve_3=T\vb_3$.
\end{itemize}
Hence $m\dot{\vv}=m\vg+T R\ve_3$. Dividing by $m$ gives \eqref{eq:m1.2}. Air drag and wind are neglected in the
model. In the toy simulation they enter as an unknown disturbance, which is exactly what the controller must reject.

\textbf{(M1.3) Attitude kinematics.} Differentiate the constraint $R^\top R=I$:
\[
  \dot R^\top R+R^\top\dot R=0\quad\Longrightarrow\quad (R^\top\dot R)^\top=-R^\top\dot R .
\]
So $R^\top\dot R$ is skew-symmetric, which means it is the hat of some vector: $R^\top\dot R=\skw{\vw}$, i.e.\\
$\dot R=R\skw{\vw}$. Physically, a point fixed on the body at $\bm r^{\mathcal B}$ sits at $R\bm r^{\mathcal B}$ in the world.
Its velocity $\dot R\bm r^{\mathcal B}=R(\vw\times\bm r^{\mathcal B})$ is a rotation about the axis $\vw$, measured in the
body frame. That is exactly what a gyroscope bolted to the drone reports.

\textbf{(M1.4) Euler's rotation equation.} The angular momentum is $J\vw$ in the body frame and $\bm h=RJ\vw$ in the
world frame. Newton's law for rotation, $\dot{\bm h}=R\vtau$ (torque in world coordinates), together with \eqref{eq:m1.3}, gives
\[
  \frac{d}{dt}(RJ\vw)=\dot RJ\vw+RJ\dot{\vw}=R\big(\skw{\vw}J\vw+J\dot{\vw}\big)=R\vtau
  \ \Longrightarrow\ J\dot{\vw}+\vw\times J\vw=\vtau .
\]
The term $\vw\times J\vw$ is the \emph{gyroscopic} coupling. A body spinning about two axes at once produces a
rotation about the third axis without any torque.

\begin{toybox}[hover, tilt and gyroscopic coupling ($m=1.5$, $J=\diag(0.02,0.02,0.04)$)]
\textbf{Hover.} $R=I$, $\vw=0$, $T=mg_0=14.715$\,N, $\vtau=0$:
$\dot{\vv}=[0,0,-9.81]+\frac{14.715}{1.5}[0,0,1]=\bm 0$. The drone hangs still.

\textbf{10$^\circ$ roll with hover thrust.} $R\ve_3=[0,-\sin10^\circ,\cos10^\circ]^\top$, so
$\dot{\vv}=g_0[0,\,-\sin10^\circ,\,\cos10^\circ-1]^\top=[0,\,-1.7035,\,-0.1490]^\top$\,m/s$^2$.
It slides sideways and sinks slightly, because only $\cos10^\circ=98.5\%$ of the thrust still holds it up.

\textbf{Yaw rate.} $R=I$, $\vw=[0,0,0.5]^\top$: $\dot R=\skw{\vw}=\begin{bmatrix}0&-0.5&0\\0.5&0&0\\0&0&0\end{bmatrix}$,
so the body $x$-axis starts to turn towards $+y$ at 0.5\,rad/s.

\textbf{Gyroscopic coupling.} $\vw=[1,0,1]^\top$, $\vtau=0$: $J\vw=[0.02,0,0.04]^\top$,
$\vw\times J\vw=[0,\,-0.02,\,0]^\top$, so $\dot{\vw}=J^{-1}[0,0.02,0]^\top=[0,\,1,\,0]^\top$\,rad/s$^2$.
\end{toybox}

\begin{projbox}
These four lines are ``the real drone''. In our set-up they are solved by Gazebo for the Iris quadrotor, while
ArduPilot SITL plays the flight computer. Our ROS~2 node never integrates them. It \emph{reads} their result on
\code{/ap/v1/pose/filtered} (position $\vp$ and orientation quaternion $\to R$) and
\code{/ap/v1/twist/filtered} ($\vv$, $\vw$). Everything else in this document either predicts this motion
(the EKF), steers it (the controller), or learns it from data (Model~2).
\end{projbox}

\begin{codebox}
\intoy \code{model1\_nonlinear\_ekf\_so3/step1\_nonlinear\_dynamics.m}, \code{lib/quad\_dynamics.m}
(RK4 integration in \code{lib/rk4\_step.m}). A 0.01\,N\,m roll-torque pulse of 0.1\,s already makes the open-loop
drone drift 2.1\,m in 4\,s (\code{figures/step1\_open\_loop.png}).
\end{codebox}

% ---------------------------------------------------------------------
\subsection{Block 2 --- Linearization: Taylor expansion and Jacobians}
\label{sec:m1b2}
% ---------------------------------------------------------------------

\begin{eqbox}{cModel}{Equations M1.5 -- M1.6 \quad First-order Taylor expansion and Jacobians}
\begin{align}
  \vf(\vx,\vu) &\approx \vf(\hx,\hu) + A\,\delta\vx + B\,\delta\vu,\qquad \delta\vx\ \text{from \eqref{eq:dx}},\ \ \delta\vu=\vu-\hu \tag{M1.5}\label{eq:m1.5}\\
  A &= \left.\frac{\partial\vf}{\partial\vx}\right|_{\hx,\hu}\in\R^{12\times12},\qquad
  B = \left.\frac{\partial\vf}{\partial\vu}\right|_{\hx,\hu}\in\R^{12\times4} \tag{M1.6}\label{eq:m1.6}
\end{align}
\end{eqbox}

\begin{dimtable}
\hx,\ \hu & operating point (where we linearize) & \R^{18},\ \R^4 & mixed\\
\delta\vx & state error $[\vp-\hp;\ \vv-\hv;\ \Log(\hR^\top R);\ \vw-\hw]$ & \R^{12} & m, m/s, rad, rad/s\\
\delta\vu=\vu-\hu & input error $[\delta T;\delta\vtau]$ & \R^4 & N, N\,m\\
A & state Jacobian & \R^{12\times12} & mixed (1/s, m/s$^2$\dots)\\
B & input Jacobian & \R^{12\times4} & mixed (1/kg, 1/(kg\,m$^2$))\\
\end{dimtable}

\subsubsection*{Derivation of the Taylor expansion}

For a smooth vector function, Taylor's theorem in several variables states
\[
  \vf(\hx+\delta\vx,\hu+\delta\vu)=\vf(\hx,\hu)+\frac{\partial\vf}{\partial\vx}\delta\vx+\frac{\partial\vf}{\partial\vu}\delta\vu
  +\mathcal O\big(\|\delta\vx\|^2+\|\delta\vu\|^2\big).
\]
Dropping the second-order remainder gives \eqref{eq:m1.5}. Entry $(i,j)$ of $A$ is $\partial f_i/\partial x_j$, i.e.\\
how much the rate of state $i$ changes when state $j$ is nudged. Because $R$ lives on $\SO$, ``nudging'' the
attitude means $R=\hR\,\Exp(\delta\vth)$, so the attitude part of $\delta\vx$ has 3 entries, not 9.

\subsubsection*{Derivation of $A$ and $B$ for the drone}

Insert $\vp=\hp+\delta\vp$, $\vv=\hv+\delta\vv$, $R=\hR\Exp(\delta\vth)\approx\hR(I+\skw{\delta\vth})$, $\vw=\hw+\delta\vw$,
$T=\hat T+\delta T$, $\vtau=\hat{\vtau}+\delta\vtau$ into \eqref{eq:m1.1}--\eqref{eq:m1.4} and keep first-order terms.

\begin{itemize}
\item \textbf{Position:} $\delta\dot{\vp}=\delta\vv$.
\item \textbf{Velocity:}
  $\dot{\vv}\approx\vg+\frac{\hat T+\delta T}{m}\hR(I+\skw{\delta\vth})\ve_3$. Using
  $\skw{\delta\vth}\ve_3=\delta\vth\times\ve_3=-\skw{\ve_3}\delta\vth$ and dropping the product $\delta T\,\delta\vth$:
  \[\delta\dot{\vv}=-\frac{\hat T}{m}\hR\skw{\ve_3}\,\delta\vth+\frac{1}{m}\hR\ve_3\,\delta T .\]
\item \textbf{Attitude:} differentiate $R=\hR\Exp(\delta\vth)$ with $\dot R=R\skw{\vw}$ and $\dot{\hR}=\hR\skw{\hw}$. To first order,
  $\skw{\delta\dot{\vth}}=\skw{\delta\vw}+\skw{\delta\vth}\skw{\hw}-\skw{\hw}\skw{\delta\vth}=\skw{\delta\vw}-\skw{\hw\times\delta\vth}$, so
  \[\delta\dot{\vth}=-\skw{\hw}\,\delta\vth+\delta\vw .\]
\item \textbf{Angular rate:} $\frac{\partial}{\partial\vw}(\vw\times J\vw)=\skw{\vw}J-\skw{J\vw}$, hence
  \[\delta\dot{\vw}=J^{-1}\big(\skw{J\hw}-\skw{\hw}J\big)\delta\vw+J^{-1}\delta\vtau .\]
\end{itemize}

Stacking the four rows gives
\begin{equation}
\label{eq:AB}
A=\begin{bmatrix}
0&I&0&0\\
0&0&-\frac{\hat T}{m}\hR\skw{\ve_3}&0\\
0&0&-\skw{\hw}&I\\
0&0&0&J^{-1}\big(\skw{J\hw}-\skw{\hw}J\big)
\end{bmatrix},\qquad
B=\begin{bmatrix}0&0\\ \frac1m\hR\ve_3&0\\ 0&0\\ 0&J^{-1}\end{bmatrix}.
\end{equation}

At hover ($\hR=I$, $\hw=0$, $\hat T=mg_0$) the key blocks become
\[
-g_0\skw{\ve_3}=\begin{bmatrix}0&9.81&0\\-9.81&0&0\\0&0&0\end{bmatrix},\qquad
\tfrac1m\ve_3=\begin{bmatrix}0\\0\\0.667\end{bmatrix},\qquad
J^{-1}=\diag(50,50,25).
\]
Read row by row: a pitch error $\delta\theta_y$ accelerates the drone along $+x$ at $9.81\,\delta\theta_y$\,m/s$^2$.
A roll error $\delta\theta_x$ accelerates it along $-y$. Extra thrust only changes vertical acceleration. Torque changes
angular acceleration with gain $1/J$.

\begin{toybox}[how good is ``first order''?]
Pitch the hovering drone by $\theta$. The linear model predicts $\dot v_x=g_0\theta$; the truth is $g_0\sin\theta$.
\begin{center}\small
\begin{tabular}{rrrr}
\toprule
$\theta$ [deg] & linear $g_0\theta$ [m/s$^2$] & true $g_0\sin\theta$ [m/s$^2$] & error\\
\midrule
1.00 & 0.1712 & 0.1712 & 0.005\%\\
2.87 (0.05 rad) & 0.4905 & 0.4903 & 0.04\%\\
10.00 & 1.7122 & 1.7035 & 0.51\%\\
20.00 & 3.4243 & 3.3552 & 2.06\%\\
28.65 (0.5 rad) & 4.9054 & 4.7035 & 4.29\%\\
\bottomrule
\end{tabular}
\end{center}
The analytic Jacobians \eqref{eq:AB} were also checked against numerical differentiation of the true flow at a
tilted, spinning, accelerating state. The largest difference is $2.3\times10^{-8}$ for $A$ and $8.4\times10^{-7}$ for $B$.
\end{toybox}

\begin{projbox}
Linearization is the ``zoom-in'' trick. Near any flight condition the complicated drone behaves like a simple
linear system, and linear systems are easy to reason about. Our EKF uses $A$ to answer one question every 10\,ms:
\emph{how fast does my uncertainty grow if I trust only the physics?} The table shows the trick is excellent for
the gentle tilts of normal flight (below 10$^\circ$, under 0.5\% error). That is why both the EKF and Model~2
can rely on it.
\end{projbox}

\begin{codebox}
\intoy \code{step2\_taylor\_jacobian.m}, \code{lib/quad\_jacobians.m}; figure \code{step2\_taylor\_accuracy.png}.
\end{codebox}

% ---------------------------------------------------------------------
\subsection{Block 3 --- Local error dynamics}
\label{sec:m1b3}
% ---------------------------------------------------------------------

\begin{eqbox}{cModel}{Equations M1.7 -- M1.9 \quad Linear dynamics of the estimation error}
\begin{align}
  \delta\vx &= \big[\vp-\hp;\ \vv-\hv;\ \Log(\hR^\top R);\ \vw-\hw\big] \tag{M1.7}\label{eq:m1.7}\\
  \delta\dot{\vx} &= A\,\delta\vx + B\,\delta\vu \tag{M1.8}\label{eq:m1.8}\\
  \delta\vx_{k+1} &= F\,\delta\vx_k + G\,\delta\vu_k,\qquad F=e^{A\Delta t}\approx I+A\Delta t,\quad G\approx B\Delta t \tag{M1.9}\label{eq:m1.9}
\end{align}
\end{eqbox}

\begin{dimtable}
\delta\vx & error between true and estimated (or nominal) state & \R^{12} & mixed\\
F & discrete error-transition matrix & \R^{12\times12} & --\\
G & discrete input matrix & \R^{12\times4} & mixed\\
\Delta t & sample time & \R & s (0.01)\\
\end{dimtable}

\subsubsection*{Derivation}

The true state obeys $\dot{\vx}=\vf(\vx,\vu)$ and the reference (estimated) state obeys the same model,
$\dot{\hx}=\vf(\hx,\hu)$. Subtracting, and using \eqref{eq:m1.5},
\[
  \delta\dot{\vx}=\vf(\vx,\vu)-\vf(\hx,\hu)\approx\big[\vf(\hx,\hu)+A\delta\vx+B\delta\vu\big]-\vf(\hx,\hu)=A\delta\vx+B\delta\vu .
\]
The large nonlinear parts cancel and only the small, linear error remains. Over one sample, with $A$ and $\delta\vu$
held constant, the solution of this linear ODE is
$\delta\vx(t_k+\Delta t)=e^{A\Delta t}\delta\vx(t_k)+\int_0^{\Delta t}e^{A s}\,ds\,B\,\delta\vu_k$.
The series $e^{A\Delta t}=I+A\Delta t+\frac12(A\Delta t)^2+\dots$ truncated after the linear term gives
\eqref{eq:m1.9}. This is exactly the forward-Euler step.

\textbf{Why the error grows without feedback.} At hover every eigenvalue of $A$ is zero ($A$ is nilpotent).
The chain $\delta\vw\to\delta\vth\to\delta\vv\to\delta\vp$ is four integrators in series, so a constant rate error
grows into a position error proportional to $t^3$. The drone is \emph{open-loop unstable}. That is why it needs both a
good estimate (block~4) and feedback (blocks 5--7).

\begin{toybox}[propagate an error for 2 s]
Extra thrust $\delta T=0.5$\,N at hover. Only row 6 of $B$ is non-zero, so
$\delta v_z(t)=\frac{\delta T}{m}t=\frac{0.5}{1.5}\cdot2=0.667$\,m/s. The linear model gives 0.6667 and the nonlinear
truth 0.6543.

Initial error $2^\circ$ roll, 10\,cm, 0.2\,rad/s yaw rate: after 2\,s the linear prediction differs from the truth by
2.6\%. With a $30^\circ$ initial roll the difference is 26\%. The linear model is \emph{local}.

With $\Delta t=0.01$\,s, $\max|F-e^{A\Delta t}|=4.9\times10^{-4}$, the size of the neglected
$\tfrac12(A\Delta t)^2$ term.
\end{toybox}

\begin{projbox}
The EKF never tries to describe the whole flight with a linear model. It only describes the
\emph{small difference} between what it believes and what is true, and that difference \emph{is} small, because the
camera and IMU keep correcting it. This is the reason for the name \emph{error-state} filter.
\end{projbox}

\begin{codebox}
\intoy \code{step3\_local\_error\_dynamics.m}; figure \code{step3\_error\_dynamics.png}.
\end{codebox}

% ---------------------------------------------------------------------
\subsection{Block 4 --- State estimation: error-state EKF}
\label{sec:m1b4}
% ---------------------------------------------------------------------

\begin{keybox}
\textbf{The idea before the maths.} The filter repeats two steps forever.
\begin{enumerate}[leftmargin=*]
  \item \textbf{Predict.} ``I know how I pushed the drone, so the physics tells me where it should be now.''
        The guess moves forward, and because physics is never perfect, the uncertainty \emph{grows}.
  \item \textbf{Correct.} ``A sensor reading has arrived. How surprised am I?'' The guess is pulled towards the
        measurement, by an amount that depends on whom we trust more, and the uncertainty \emph{shrinks}.
\end{enumerate}
The filter stores two things: the \textbf{estimate} $\hx_k=(\hp,\hv,\hR,\hw)$ (its best guess) and the
\textbf{covariance} $P_k$ (how unsure it is about each of the 12 error directions of \eqref{eq:dx}).
\end{keybox}

\textbf{Assumptions about noise.} The real drone follows the model only approximately. Each step it is pushed by
a small unknown random error $\bm w_k$ (wind, model mistakes) with covariance $Q$. Every sensor reads the true
value plus a random error $\bm v_k$ with covariance $R_n$:
\[
  \vx_{k+1}=\vf_d(\vx_k,\vu_k)\ \text{plus a random error }\bm w_k\sim\mathcal N(0,Q),\qquad
  \vz_k=h(\vx_k)+\bm v_k,\quad \bm v_k\sim\mathcal N(0,R_n).
\]
Here $\vf_d$ is one simulation step (RK4) of \eqref{eq:m1.1}--\eqref{eq:m1.4}, and $\mathcal N(0,Q)$ means ``Gaussian
with mean zero and covariance $Q$''. The errors are added in the 12 error directions, so for the attitude
``adding'' means a small rotation, as in \eqref{eq:correct}.

\begin{eqbox}{cEst}{Equations M1.10 -- M1.14 \quad Error-state EKF}
\textbf{Time update (prediction)}
\begin{align}
  \hx_{k+1} &= \vf_d(\hx_k,\vu_k) \tag{M1.10}\label{eq:m1.10}\\
  P_{k+1}   &= F P_k F^\top + Q \tag{M1.11}\label{eq:m1.11}
\end{align}
\textbf{Measurement update (correction)}
\begin{align}
  K   &= P H^\top\big(HPH^\top+R_n\big)^{-1} \tag{M1.12}\label{eq:m1.12}\\
  \delta\hat{\vx} &= K\,\vr,\qquad \vr=\vz-h(\hx)\ \ \text{(innovation)},\qquad
  \text{then correct } \hx \text{ by } \delta\hat{\vx} \text{ using \eqref{eq:correct}} \tag{M1.13}\label{eq:m1.13}\\
  P   &\leftarrow (I-KH)\,P \tag{M1.14}\label{eq:m1.14}
\end{align}
Written out, the correction in (M1.13) is
$\hp\leftarrow\hp+\delta\hat{\vp}$, $\hv\leftarrow\hv+\delta\hat{\vv}$, $\hR\leftarrow\hR\,\Exp(\delta\hat{\vth})$,
$\hw\leftarrow\hw+\delta\hat{\vw}$.
\end{eqbox}

\begin{dimtable}
\hx_k & estimated state $(\hp,\hv,\hR,\hw)$ & \R^{18} & mixed\\
P_k & error covariance & \R^{12\times12} & m$^2$, (m/s)$^2$, rad$^2$, \dots\\
Q & process-noise covariance per step (wind, model error) & \R^{12\times12} & as $P$\\
\vz & measurement (camera pose: $m=6$, gyro: $m=3$) & \R^m & m, rad, rad/s\\
h(\cdot) & measurement model & \R^{18}\to\R^m & --\\
H=\partial h/\partial\delta\vx & measurement Jacobian & \R^{m\times12} & --\\
R_n & measurement-noise covariance & \R^{m\times m} & m$^2$, rad$^2$, \dots\\
K & Kalman gain & \R^{12\times m} & mixed\\
\vr=\vz-h(\hx) & innovation: measured minus expected (``surprise'') & \R^m & as $\vz$\\
\delta\hat{\vx} & correction computed from the innovation & \R^{12} & mixed\\
\end{dimtable}

\subsubsection*{Derivation}

\textbf{(M1.10) Mean prediction.} The truth is the estimate corrected by an unknown error $\delta\vx_k$ whose
average is zero (a good estimate is not biased to one side). By the linearization of block~3, one step later the
truth is $\vf_d(\hx_k,\vu_k)$ corrected by $F\delta\vx_k$. Since $\delta\vx_k$ averages to zero, so does $F\delta\vx_k$,
and the best guess of the next state is simply $\vf_d(\hx_k,\vu_k)$. The full nonlinear model is used for the guess
itself. The linear model is only needed for the uncertainty.

\textbf{(M1.11) Covariance prediction.} From \eqref{eq:m1.9} plus noise, $\delta\vx_{k+1}=F\delta\vx_k+\bm w_k$. Then
\[
  P_{k+1}=\mathbb E[\delta\vx_{k+1}\delta\vx_{k+1}^\top]
  =F\,\mathbb E[\delta\vx_k\delta\vx_k^\top]F^\top+\mathbb E[\bm w_k\bm w_k^\top]
  =FP_kF^\top+Q,
\]
because the noise is independent of the current error, so the cross terms vanish. Uncertainty is stretched by
the dynamics ($F\cdot F^\top$) and inflated by the noise ($+Q$).

\textbf{(M1.12) Kalman gain.} Linearize the measurement: the innovation $\vr=\vz-h(\hx)\approx H\delta\vx+\bm v$
(``what I see minus what I expected = my error seen through the sensor, plus sensor noise'').
We correct the error estimate by a linear rule $\delta\hat{\vx}=K\vr$. The error after the correction is
$\delta\vx^+=\delta\vx-K(H\delta\vx+\bm v)=(I-KH)\delta\vx-K\bm v$, whose covariance is
\begin{equation}
\label{eq:joseph}
  P^+=(I-KH)P(I-KH)^\top+KR_nK^\top \qquad\text{(Joseph form, valid for any $K$).}
\end{equation}
The best $K$ minimises the total variance $\tr P^+$. Using $\frac{\partial}{\partial K}\tr(KMK^\top)=2KM$ for
symmetric $M$ and $\frac{\partial}{\partial K}\tr(KN)=N^\top$:
\[
  \frac{\partial\,\tr P^+}{\partial K}=-2PH^\top+2K\big(HPH^\top+R_n\big)=0
  \ \Longrightarrow\ K=PH^\top\underbrace{\big(HPH^\top+R_n\big)}_{S\ \text{(innovation covariance)}}{}^{-1}.
\]

\textbf{(M1.13) State update.} The best estimate of the error is $\delta\hat{\vx}=K\vr$. It is applied to the
estimate with the correction rule \eqref{eq:correct}: positions, velocities and rates are added, and the attitude
is rotated, $\hR\leftarrow\hR\Exp(\delta\hat{\vth})$, so $\hR$ stays a valid rotation matrix. The error has now been
absorbed into the estimate, so the filter starts the next step believing its error is zero on average again.

\textbf{(M1.14) Covariance update.} Expand \eqref{eq:joseph} and use $KS=PH^\top$:
$P^+=P-KHP-PH^\top K^\top+KSK^\top=P-KHP-PH^\top K^\top+PH^\top K^\top=(I-KH)P$.
The toy code uses this short form and re-symmetrises $P$ after every step.

\subsubsection*{Our measurement models}

\begin{center}\small
\begin{tabular}{@{}llll@{}}
\toprule
Sensor (rate) & innovation $\vr$ & $H$ & noise $\sqrt{\diag R_n}$\\
\midrule
Gyroscope (100\,Hz) & $\tilde{\vw}-\hw$ & $[\,0\ \ 0\ \ 0\ \ I\,]$ & 0.01\,rad/s\\
Camera / VIO pose (25\,Hz) & $\begin{bmatrix}\tilde{\vp}-\hp\\ \Log(\hR^\top\tilde R)\end{bmatrix}$ &
  $\begin{bmatrix}I&0&0&0\\0&0&I&0\end{bmatrix}$ & 0.05\,m, 0.02\,rad\\
\bottomrule
\end{tabular}
\end{center}
Each block of $H$ is $3\times3$, and the columns follow $[\delta\vp,\delta\vv,\delta\vth,\delta\vw]$. Velocity is
\emph{never measured}. The filter infers it from how the measured position changes and from the physics in $F$.

\begin{toybox}[one predict + one correct, altitude only]
Prior: altitude $\hat z_k=10$\,m, variance $P_k=3.5$\,m$^2$. Commanded climb 0.5\,m/s for $\Delta t=1$\,s, $F=1$, $Q=0.5$:
\[\hat z^-=10+0.5=10.5\ \text{m},\qquad P^-=1\cdot3.5\cdot1+0.5=4.0\ \text{m}^2.\]
GPS reads $z=12.5$\,m with $R_n=1$\,m$^2$, $H=1$:
\[K=\frac{4}{4+1}=0.8,\quad \hat z^+=10.5+0.8\,(12.5-10.5)=12.1\ \text{m},\quad P^+=(1-0.8)\cdot4=0.8\ \text{m}^2.\]
The GPS is four times more trustworthy than the prediction ($P^-/R_n=4$), so the estimate moves 80\% of the way
towards it, and the variance shrinks from 4.0 to 0.8.

\textbf{Full 12-state filter} (20\,s helix flight, gusts, noisy sensors, wrong initial guess):
\begin{center}\small
\begin{tabular}{lrl}
\toprule
quantity & RMSE after 2\,s & for comparison\\
\midrule
position & 2.55\,cm & raw camera: 8.55\,cm\\
velocity (never measured) & 3.9\,cm/s & --\\
attitude & 0.20$^\circ$ & camera noise: 1.15$^\circ$\\
body rate & 0.0025\,rad/s & gyro noise: 0.01\,rad/s\\
\bottomrule
\end{tabular}
\end{center}
In 99\% of the samples all 12 errors lie inside the $\pm3\sigma$ band predicted by $P_k$, so the filter is
\emph{consistent}: it knows how uncertain it is.
\end{toybox}

\begin{figure}[h]
  \centering
  \includegraphics[width=0.49\linewidth]{step4_ekf_trajectory.png}\hfill
  \includegraphics[width=0.49\linewidth]{step4_ekf_errors.png}
  \caption{Toy model, Model~1 step~4. Left: noisy camera positions (grey), truth (black) and EKF estimate.
  Right: estimation errors stay inside the $\pm3\sigma$ band computed from $P_k$.}
  \label{fig:m1-ekf}
\end{figure}

\begin{projbox}[this is the ``VIO'' of our title]
A monocular camera alone cannot tell scale (a small room close up looks like a big room far away), and an IMU
alone drifts within seconds. The EKF is the referee that fuses them. Between camera frames it predicts with the
physics and the gyro. When a frame arrives, the visual-odometry pose corrects the prediction, and $K$ decides how much
to trust each source. Its output $\hx$ (position, velocity, attitude, rate) is the only thing the controller ever
sees. In the current SITL set-up ArduPilot's own EKF3 produces \code{/ap/v1/*/filtered}. Our VIO pipeline supplies
the camera measurement $\vz_{\text{cam}}$ to exactly this structure.
\end{projbox}

\begin{codebox}
\intoy \code{step4\_error\_state\_ekf.m}, \code{lib/ekf\_predict.m}, \code{lib/ekf\_update.m}.
\end{codebox}

% ---------------------------------------------------------------------
\subsection{Block 5 --- Position control (outer loop)}
\label{sec:m1b5}
% ---------------------------------------------------------------------

\begin{eqbox}{cOuter}{Equations M1.15 -- M1.16 \quad Tracking errors and desired acceleration}
\begin{align}
  \ve_p=\vp_d-\hp,\qquad \ve_v&=\vv_d-\hv \tag{M1.15}\label{eq:m1.15}\\
  \va_d&=\ddot{\vp}_d+K_p\,\ve_p+K_v\,\ve_v \tag{M1.16}\label{eq:m1.16}
\end{align}
\end{eqbox}

\begin{dimtable}
\vp_d,\ \vv_d,\ \ddot{\vp}_d & desired position, velocity, acceleration & \R^3 & m, m/s, m/s$^2$\\
\ve_p,\ \ve_v & position and velocity error & \R^3 & m, m/s\\
K_p & position gain (diagonal, positive) & \R^{3\times3} & 1/s$^2$\\
K_v & velocity gain (diagonal, positive) & \R^{3\times3} & 1/s\\
\va_d & desired acceleration & \R^3 & m/s$^2$\\
\end{dimtable}

\subsubsection*{Derivation}

Suppose the inner loops achieve the commanded acceleration exactly, $\ddot{\vp}=\va_d$. Then the position error
obeys
\[
  \ddot{\ve}_p=\ddot{\vp}_d-\ddot{\vp}=\ddot{\vp}_d-\va_d=-K_p\ve_p-K_v\dot{\ve}_p
  \quad\Longrightarrow\quad \ddot{\ve}_p+K_v\dot{\ve}_p+K_p\ve_p=\bm0 .
\]
With diagonal gains, each axis is a mass--spring--damper with characteristic polynomial $s^2+k_vs+k_p$.
Both roots have negative real part if and only if $k_p,k_v>0$, so the error decays to zero. Matching with
$s^2+2\zeta\omega_ns+\omega_n^2$:
\[
  \omega_n=\sqrt{k_p},\qquad \zeta=\frac{k_v}{2\sqrt{k_p}},\qquad t_{2\%}\approx\frac{4}{\zeta\omega_n}.
\]
The feed-forward $\ddot{\vp}_d$ supplies the acceleration a perfect pilot would need. The feedback terms only fix what
is left over.

\begin{center}\small
\begin{tabular}{@{}lrrrrr@{}}
\toprule
axis & $k_p$ & $k_v$ & $\omega_n$ [rad/s] & $\zeta$ & $t_{2\%}$ [s]\\
\midrule
$x,\ y$ & 1.5 & 1.2 & 1.225 & 0.490 & 6.7\\
$z$ & 2.0 & 1.5 & 1.414 & 0.530 & 5.3\\
\bottomrule
\end{tabular}
\end{center}
These are our \code{controller.yaml} gains: gentle, slightly under-damped, a little overshoot, settling in about 6\,s.

\begin{toybox}[50 cm low and climbing]
Setpoint $\vp_d=[0,0,2]$\,m (the \code{controller.yaml} hover point), $\vv_d=\ddot{\vp}_d=\bm0$.
Estimate $\hp=[0.1,\,-0.2,\,1.5]$\,m, $\hv=[0,0,0.3]$\,m/s.
\[
\ve_p=[-0.1,\ 0.2,\ 0.5],\quad \ve_v=[0,\ 0,\ -0.3],\quad
\va_d=\begin{bmatrix}1.5(-0.1)\\1.5(0.2)\\2(0.5)+1.5(-0.3)\end{bmatrix}=\begin{bmatrix}-0.15\\0.30\\0.55\end{bmatrix}\text{ m/s}^2 .
\]
It is 50\,cm low, so it wants to accelerate up. It is already climbing at 0.3\,m/s, so it asks for less (0.55 instead of 1.0).
\end{toybox}

\begin{projbox}
This is the ``where am I versus where should I be'' layer: \code{SO3Controller.compute\_position\_errors()} and the
first line of \code{compute\_desired\_force()} in \code{so3\_controller.py}. The positions come from the EKF, and the setpoint
comes from \code{desired\_x/y/z} in the YAML (or a waypoint generator). Large $K_p$ means an aggressive drone, and large
$K_v$ means a damped, smooth drone.
\end{projbox}

\begin{codebox}
\intoy \code{step5\_position\_control.m}, \code{lib/position\_control.m}; figure \code{step5\_error\_decay.png}.
\end{codebox}

% ---------------------------------------------------------------------
\subsection{Block 6 --- Force and attitude generation}
\label{sec:m1b6}
% ---------------------------------------------------------------------

\begin{eqbox}{cGuid}{Equations M1.17 -- M1.19 \quad Desired force, thrust direction and attitude}
\begin{align}
  \vF_d &= m\,(\va_d-\vg) \tag{M1.17}\label{eq:m1.17}\\
  \vb_{3d} &= \frac{\vF_d}{\|\vF_d\|} \tag{M1.18}\label{eq:m1.18}\\
  \vb_{2d}=\frac{\vb_{3d}\times\vb_{1c}}{\|\vb_{3d}\times\vb_{1c}\|},\quad
  \vb_{1d}&=\vb_{2d}\times\vb_{3d},\quad
  R_d=\big[\vb_{1d}\ \ \vb_{2d}\ \ \vb_{3d}\big],\quad \vb_{1c}=[\cos\psi_d,\ \sin\psi_d,\ 0]^\top \tag{M1.19}\label{eq:m1.19}
\end{align}
\end{eqbox}

\begin{dimtable}
\vF_d & desired total force vector (world) & \R^3 & N\\
\vb_{1d},\vb_{2d},\vb_{3d} & desired body axes (unit vectors) & \R^3 & --\\
\psi_d & desired yaw (heading) & \R & rad\\
\vb_{1c} & heading direction in the horizontal plane & \R^3 & --\\
R_d & desired attitude & \SO & --\\
\end{dimtable}

\subsubsection*{Derivation}

\textbf{(M1.17)} We want $\dot{\vv}=\va_d$. From \eqref{eq:m1.2},
$\va_d=\vg+\frac{T}{m}R\ve_3$, so the thrust vector must satisfy $T\,R\ve_3=m(\va_d-\vg)=\vF_d$.
In ENU $-\vg=[0,0,+g_0]$, so even for $\va_d=0$ the drone must push up with $mg_0$. In the code's notation,
$\vF_d=m(\va_{\text{cmd}}+g\ve_3)$, which is the same expression.

\textbf{(M1.18) Underactuation.} A quadrotor has 4 inputs but 6 degrees of freedom. It cannot push sideways. The
only way to produce $\vF_d$ is to point the body axis $\vb_3=R\ve_3$ along it. Because $\|R\ve_3\|=1$, matching
$T\vb_3=\vF_d$ forces $T=\|\vF_d\|$ and $\vb_{3d}=\vF_d/\|\vF_d\|$. Translation is controlled \emph{through} the attitude.

\textbf{(M1.19) Completing the rotation.} $\vb_{3d}$ fixes two of the three attitude angles. The remaining freedom is
yaw, which we choose. By construction:
$\vb_{2d}\perp\vb_{3d}$ and $\vb_{2d}\perp\vb_{1c}$ (cross product); $\vb_{1d}=\vb_{2d}\times\vb_{3d}$ is perpendicular to both
and has unit length, since it is the cross product of orthogonal unit vectors. The columns are therefore orthonormal, and since
$\vb_{1d}\times\vb_{2d}=\vb_{3d}$ the frame is right-handed. Hence $R_d^\top R_d=I$ and $\det R_d=+1$, so $R_d\in\SO$.
$\vb_{1d}$ is the projection of the desired heading onto the plane perpendicular to the thrust.
The construction fails only if $\vb_{3d}\parallel\vb_{1c}$ (a 90$^\circ$ tilt). The code then nudges $\psi_d$ by
$10^{-3}$\,rad.

\begin{toybox}[continuing block 5]
$\va_d=[-0.15,0.30,0.55]$, $m=1.5$, $\vg=[0,0,-9.81]$, $\psi_d=0$:
\[
\vF_d=1.5\,[-0.15,\ 0.30,\ 10.36]=[-0.225,\ 0.450,\ 15.540]\ \text{N},\qquad \|\vF_d\|=15.548\ \text{N},
\]
\[
\vb_{3d}=[-0.0145,\ 0.0289,\ 0.9995],\qquad
R_d=\begin{bmatrix}0.9999&0.0000&-0.0145\\0.0004&0.9996&0.0289\\0.0145&-0.0289&0.9995\end{bmatrix}.
\]
Tilt $=\arccos(0.9995)=1.85^\circ$: lean slightly towards $-x$, $+y$ and push 15.5\,N (more than the 14.7\,N weight,
because it must also climb). Check: $\|R_d^\top R_d-I\|=1.1\times10^{-16}$, $\det R_d=1$.
\end{toybox}

\begin{projbox}
A drone moves sideways by \emph{leaning}, like a person pushing a heavy cart. This block converts ``accelerate
this way'' into ``lean this much in this direction, and push this hard''. It is \code{compute\_desired\_force()}
and \code{compute\_desired\_attitude()} in \code{so3\_controller.py}. The attitude is built as a rotation matrix,
never as roll/pitch/yaw angles, so there is no gimbal lock.
\end{projbox}

\begin{codebox}
\intoy \code{step6\_force\_attitude.m}, \code{lib/force\_attitude.m}.
\end{codebox}

% ---------------------------------------------------------------------
\subsection{Block 7 --- Attitude control on SO(3) (inner loop)}
\label{sec:m1b7}
% ---------------------------------------------------------------------

\begin{eqbox}{cInner}{Equations M1.20 -- M1.23 \quad Geometric attitude controller}
\begin{align}
  \ve_R &= \tfrac12\big(R_d^\top\hR-\hR^\top R_d\big)^\vee \tag{M1.20}\label{eq:m1.20}\\
  \ve_\omega &= \hw-\hR^\top R_d\,\vw_d \tag{M1.21}\label{eq:m1.21}\\
  \vtau &= -K_R\ve_R-K_\omega\ve_\omega+\hw\times J\hw
          -J\big(\skw{\hw}\hR^\top R_d\vw_d-\hR^\top R_d\dot{\vw}_d\big) \tag{M1.22}\label{eq:m1.22}\\
  T &= \vF_d\cdot\hR\,\ve_3 \tag{M1.23}\label{eq:m1.23}
\end{align}
\end{eqbox}

\begin{dimtable}
\ve_R & attitude error (``$\approx$ rotation vector'' for small errors) & \R^3 & rad (small angles)\\
\ve_\omega & angular-velocity error (body frame) & \R^3 & rad/s\\
\vw_d,\ \dot{\vw}_d & desired body rate and its derivative & \R^3 & rad/s, rad/s$^2$\\
K_R,\ K_\omega & attitude and rate gains & \R^{3\times3} & N\,m, N\,m\,s\\
\vtau,\ T & torque and thrust commands & \R^3,\ \R & N\,m, N\\
\end{dimtable}

\subsubsection*{Derivation of $\ve_R$ (M1.20)}

Measure ``how far'' $\hR$ is from $R_d$ with the error function
\[
  \Psi(\hR,R_d)=\tfrac12\tr\big(I-R_d^\top\hR\big)=1-\cos\varphi\ \ \in[0,2],
\]
where $\varphi$ is the angle of the relative rotation $R_d^\top\hR$ (because $\tr R=1+2\cos\varphi$). $\Psi=0$
exactly when $\hR=R_d$. Now rotate the drone slightly, $\hR\to\hR\Exp(\epsilon\bm\eta)$, and differentiate:
\[
  \left.\frac{d}{d\epsilon}\Psi\right|_{\epsilon=0}=-\tfrac12\tr\big(R_d^\top\hR\skw{\bm\eta}\big)
  \overset{\eqref{eq:id3}}{=}\tfrac12\bm\eta^\top\big(R_d^\top\hR-\hR^\top R_d\big)^\vee=\ve_R\cdot\bm\eta .
\]
So $\ve_R$ is the \emph{gradient} of the attitude error on $\SO$: it points in the direction of rotation that increases
the error fastest, and the controller pushes the opposite way. For a single-axis error of angle $\varphi$ about $\bm n$,
$\ve_R=\sin\varphi\,\bm n$, which is $\approx\varphi\bm n$ for small angles.

\subsubsection*{Derivation of $\ve_\omega$ (M1.21)}

Let $R_e=R_d^\top\hR$ be the relative rotation. Using $\dot{\hR}=\hR\skw{\hw}$, $\dot R_d=R_d\skw{\vw_d}$ and
$\skw{\vw_d}R_e=R_e\skw{R_e^\top\vw_d}$ (from \eqref{eq:id2}),
\[
  \dot R_e=R_d^\top\hR\skw{\hw}-\skw{\vw_d}R_d^\top\hR=R_e\,\skw{\hw-\hR^\top R_d\vw_d}=R_e\,\skw{\ve_\omega}.
\]
$\ve_\omega$ is therefore the angular velocity of the \emph{error} rotation, expressed in the body frame: how fast the drone
is turning relative to where it should be turning. Note the sign: \emph{actual minus desired}.

\subsubsection*{Derivation of the torque law (M1.22) by Lyapunov's method}

Differentiate $\ve_\omega$. With $\frac{d}{dt}(\hR^\top R_d)=-\skw{\hw}\hR^\top R_d+\hR^\top R_d\skw{\vw_d}$ and
$\skw{\vw_d}\vw_d=0$,
\[
  J\dot{\ve}_\omega=J\dot{\hw}+J\skw{\hw}\hR^\top R_d\vw_d-J\hR^\top R_d\dot{\vw}_d
  \overset{\eqref{eq:m1.4}}{=}\vtau-\hw\times J\hw+J\big(\skw{\hw}\hR^\top R_d\vw_d-\hR^\top R_d\dot{\vw}_d\big).
\]
Substituting \eqref{eq:m1.22}, every nonlinear term cancels and the error dynamics become
\[
  J\dot{\ve}_\omega=-K_R\ve_R-K_\omega\ve_\omega .
\]
Take $K_R=k_RI$, $K_\omega=k_\omega I$ and the energy-like function
$V=\tfrac12\ve_\omega^\top J\ve_\omega+k_R\Psi$. Using $\dot\Psi=\ve_R\cdot\ve_\omega$ (same calculation as above,
with $\bm\eta=\ve_\omega$):
\[
  \dot V=\ve_\omega^\top\big(-k_R\ve_R-k_\omega\ve_\omega\big)+k_R\,\ve_R\cdot\ve_\omega=-k_\omega\|\ve_\omega\|^2\le0 .
\]
$V$ never increases. LaSalle's invariance principle then shows that $(\ve_R,\ve_\omega)\to0$ from every attitude with
$\Psi<2$, i.e.\ from everything except exactly upside-down. Unlike Euler-angle controllers, there is no singularity.

\textbf{Simplified form.} Near the target ($\hR^\top R_d\approx I$) and for slow manoeuvres, the last bracket reduces to
$-J(\skw{\hw}\vw_d-\dot{\vw}_d)\approx+J\dot{\vw}_d$, which is the term printed in the original image.

\subsubsection*{Derivation of the thrust (M1.23)}

Thrust acts only along the \emph{current} axis $\hR\ve_3$. Its useful part is the projection of the desired force on
that axis, $T=\vF_d\cdot\hR\ve_3$. When the attitude has converged ($\hR\ve_3=\vb_{3d}$), $T=\|\vF_d\|$. While still
tilted the wrong way, $T$ is automatically reduced, so the drone does not push hard in a wrong direction.

\begin{toybox}[recovering from a 10$^\circ$ roll]
$\hR=\Exp([0.1745,0,0])$ (10$^\circ$ roll), $\hw=[0.2,0,0]$\,rad/s (still rolling further), $R_d=I$,
$\vw_d=\dot{\vw}_d=0$, $\vF_d=[0,0,14.715]$:
\[
\ve_R=\tfrac12(\hR-\hR^\top)^\vee=[\sin10^\circ,0,0]=[0.1736,0,0],\qquad \ve_\omega=[0.2,0,0],
\]
\[
\vtau=-\begin{bmatrix}4(0.1736)\\0\\0\end{bmatrix}-\begin{bmatrix}0.8(0.2)\\0\\0\end{bmatrix}+\underbrace{\hw\times J\hw}_{=0}
=\begin{bmatrix}-0.8546\\0\\0\end{bmatrix}\text{N\,m},\qquad
T=14.715\cos10^\circ=14.491\ \text{N}.
\]
A restoring torque about $-x$ rolls it back, and thrust is reduced to the useful vertical part.
\end{toybox}

\begin{fixbox}[sign of the rate error]
The image prints $\ve_\omega=\vw_d-\hw$ together with $\vtau=-K_\omega\ve_\omega+\dots$, i.e.\ $\vtau=+K_\omega\hw+\dots$.
That is negative damping. With it, $\dot V=+k_\omega\|\ve_\omega\|^2$ and the energy \emph{grows}. The toy model starts
60$^\circ$ from level: with the corrected sign the error decays to $2\times10^{-6}$\,deg within 3\,s and $V$ decreases
monotonically. With the image's sign the drone exceeds 50\,rad/s (tumbling) after 0.066\,s (\cref{fig:m1-att}).
\end{fixbox}

\begin{figure}[h]
  \centering
  \includegraphics[width=0.92\linewidth]{step7_attitude_recovery.png}
  \caption{Toy model, Model~1 step~7. Left: recovery from 60$^\circ$ with the corrected and the original sign of
  $\ve_\omega$. Right: the Lyapunov function $V$ decreases monotonically with the corrected law.}
  \label{fig:m1-att}
\end{figure}

\begin{projbox}
This is the fast inner loop: ``I should be leaning like $R_d$ but I am leaning like $\hR$, so which way do I
twist?''. It works directly with rotation matrices, so it handles any orientation without gimbal lock. That is the
central claim of our README (``Euler-angle-free''). In code it is \code{compute\_attitude\_error()},
\code{compute\_angular\_velocity\_error()} and \code{compute\_moment()} in \code{so3\_controller.py}. In the current
ROS~2 interface the node publishes $\vF_d$, $\ve_R$ and $\vtau$ as diagnostics and sends velocity commands, and ArduPilot
executes its own inner loop. The equations are the same.
\end{projbox}

\begin{codebox}
\intoy \code{step7\_attitude\_control.m}, \code{lib/attitude\_control.m}.
\end{codebox}

% ---------------------------------------------------------------------
\subsection{Block 8 --- Actuator commands and feedback}
\label{sec:m1b8}
% ---------------------------------------------------------------------

\begin{eqbox}{cPlant}{Equations M1.24 -- M1.25 \quad Control allocation and motor speeds}
\begin{align}
  \begin{bmatrix}T\\ \vtau\end{bmatrix}&=\Gamma\,\vf,\qquad
  \Gamma=\begin{bmatrix}1&1&1&1\\ y_1&y_2&y_3&y_4\\ -x_1&-x_2&-x_3&-x_4\\ \kappa_1c_m&\kappa_2c_m&\kappa_3c_m&\kappa_4c_m\end{bmatrix} \tag{M1.24}\label{eq:m1.24}\\
  \Omega_i&=\sqrt{\frac{\big(\Gamma^{-1}\vu\big)_i}{k_f}}\quad\text{clipped to }[\Omega_{\min},\Omega_{\max}] \tag{M1.25}\label{eq:m1.25}
\end{align}
\end{eqbox}

\begin{dimtable}
\vf=[f_1..f_4] & individual rotor thrusts & \R^4 & N\\
(x_i,y_i) & rotor position in the body frame & \R^2 & m\\
\kappa_i & spin sign ($-1$ CCW, $+1$ CW) & \{\pm1\} & --\\
c_m & yaw torque per newton of thrust & \R & m\\
\Gamma & allocation (mixer) matrix & \R^{4\times4} & --, m\\
k_f & thrust constant & \R & N\,s$^2$\\
\vOm & rotor speeds & \R^4 & rad/s\\
\end{dimtable}

\subsubsection*{Derivation}

Each rotor $i$ produces a force $f_i\ve_3$ at $\bm r_i=[x_i,y_i,0]^\top$. Summing:
\begin{itemize}
  \item total thrust $T=\sum_if_i$;
  \item torque $\sum_i\bm r_i\times f_i\ve_3=\sum_if_i[y_i,\,-x_i,\,0]^\top$, which gives the roll and pitch rows;
  \item yaw: the air resists the spinning propeller, so the motor reacts on the body with $\kappa_ic_mf_i$ about $\ve_3$.
        CCW rotors push the body clockwise.
\end{itemize}
This gives \eqref{eq:m1.24}. Momentum theory gives a rotor thrust proportional to the square of the tip speed,
$f_i=k_f\Omega_i^2$. Inverting both relations gives \eqref{eq:m1.25}. Motors can neither reverse nor exceed their maximum
speed, so $\Omega_i$ is clipped, and the wrench actually produced is $\Gamma k_f\vOm^2$.

For the Iris (ArduPilot Quad-X: 1 front-right CCW, 2 back-left CCW, 3 front-left CW, 4 back-right CW):
\[
\Gamma=\begin{bmatrix}1&1&1&1\\-0.22&0.22&0.22&-0.22\\-0.13&0.13&-0.13&0.13\\-0.016&-0.016&0.016&0.016\end{bmatrix}.
\]

\begin{toybox}[hover and the roll correction of block 7]
\textbf{Hover} $\vu=[14.715,0,0,0]$: $f_i=14.715/4=3.679$\,N,
$\Omega_i=\sqrt{3.679/8.549\times10^{-6}}=656$\,rad/s $=6264$\,rpm.

\textbf{Roll correction} $\vu=[14.491,-0.8546,0,0]$: by symmetry $f_1=f_4=a$ and $f_2=f_3=b$, with
$2a+2b=14.491$ and $0.44(b-a)=-0.8546$. So $a=4.594$\,N and $b=2.652$\,N, i.e.\\
$\vOm=[733,\,557,\,557,\,733]$\,rad/s. The right-side rotors (1, 4) speed up and the drone rolls back to level.
\end{toybox}

\begin{figure}[h]
  \centering
  \includegraphics[width=0.43\linewidth]{step8_closed_loop_3d.png}\hfill
  \includegraphics[width=0.55\linewidth]{step8_closed_loop_signals.png}
  \caption{Toy model, Model~1 step~8: the complete loop (truth $\to$ sensors $\to$ EKF $\to$ blocks 5--7 $\to$
  allocation $\to$ truth) tracking a climbing helix in wind. After the transient: tracking RMSE 6.8\,cm,
  estimation RMSE 2.6\,cm, motor speeds 619--709\,rad/s.}
  \label{fig:m1-loop}
\end{figure}

\begin{projbox}
The flight controller does not ``send a torque''. It sends four motor speeds. This block is ArduPilot's motor
mixer. When our node commands velocities (\code{/ap/v1/cmd\_vel}), ArduPilot runs blocks 5--8 internally. When it
commands force and moment, this is the matrix that turns them into ESC signals. The closed loop in
\cref{fig:m1-loop} is the whole of Model~1 working together: the arrow from block~8 back to block~1.
\end{projbox}

\begin{codebox}
\intoy \code{step8\_actuators\_feedback.m}, \code{lib/motor\_allocation.m}.
\end{codebox}

% ---------------------------------------------------------------------
\subsection{Recap of Model 1}
\label{sec:m1recap}
% ---------------------------------------------------------------------

\begin{center}\small
\begin{tabularx}{\linewidth}{@{}c >{\raggedright\arraybackslash}p{2.9cm} >{\raggedright\arraybackslash}p{3.2cm} >{\raggedright\arraybackslash}p{3.0cm} X@{}}
\toprule
\textbf{\#} & \textbf{Block} & \textbf{In} & \textbf{Out} & \textbf{The idea in one sentence}\\
\midrule
1 & Nonlinear dynamics & thrust $T$, torque $\vtau$ & motion $\vp,\vv,R,\vw$ & Newton and Euler: push along the body $z$-axis, gravity pulls down, torque spins the body.\\
2 & Linearization & operating point & $A$, $B$ & Near any flight condition the drone behaves like a linear system.\\
3 & Error dynamics & $A$, $B$, $\Delta t$ & $F$ & Small errors evolve linearly; without feedback they grow.\\
4 & Error-state EKF & $\vu_k$, sensors $\vz_k$ & estimate $\hx_k$, $P_k$ & Predict with physics, correct with camera and IMU, weighting by trust.\\
5 & Position control & $\hp,\hv$, reference & $\va_d$ & Accelerate towards the target like a spring with a damper.\\
6 & Force \& attitude & $\va_d$, yaw & $\vF_d$, $R_d$ & To accelerate sideways, lean; thrust direction $=$ desired force direction.\\
7 & SO(3) attitude & $\hR,\hw,R_d,\vF_d$ & $\vtau$, $T$ & Twist along the gradient of the attitude error; energy always decreases.\\
8 & Allocation & $T$, $\vtau$ & $\Omega_1..\Omega_4$ & Split thrust and torque among four rotors; speed $\propto\sqrt{\text{thrust}}$.\\
\bottomrule
\end{tabularx}
\end{center}
